%

\section{Towards Theoretical Guarantees for\\ Perturbation-Based Methods}\label{sec:sample-complexity}

We have shown both theoretically and empirically that \completename{} and \additivename{} \fielddefnames{} like \shap{} and \intgradshort{} can be unreliable for solving \downtasknames{} like algorithmic recourse and spurious \covname{} identification.
At the same time, while simpler \fielddefnames{} like \vanillagrad{} can be accurate for inferring sufficiently \localname{} \modbehaviourname{}, this can also clearly fail even for simple tasks that demand understanding a moderately \localname{} region.
So, faced with solving a given \downtaskname{}, what should a practitioner do?

To answer this question, we turn to the existing literature on \fielddefnames{} that are empirically better than \shap{} and \intgradshort{} for many \tasknames{}.
For example,
\citet{fong2017interpretable} learn a blurring mask that changes the class probability while simultaneously maximizing how ``informative'' the mask is,
\citet{petsiuk18rise} use randomly generated masks to learn which perturbations maximally change the model output,
\citet{kapishnikov19xrai} use image masking to find regions that maximally change the class probability,
\citet{qi2019visualizing} and \citet{khorram2021igos++} learn a feature mask based on both removing and adding components to the image to change the class probability,
and \citet{shitole2021one} combine multiple saliency maps into a single graph to illustrate multiple different minimal perturbations to change the model output.


However, despite the many methods listed above, mathematically establishing when these methods are (and are not) reliable has remained elusive.
As a first step towards such theory, we consider a \methodname{} that is \emph{guaranteed} to work: brute-force solving \downtasknames{} via repeated \modelname{} evaluations.
In particular, since \countmodbehaviourname{} is determined by $\model(\covvaldum)$ for $\covvaldum$ near $\covval$, this can always be inferred by computing $\model(\covvaldum)$ (or, perhaps, $\grad \model(\covvaldum)$) at sufficiently many \obsnames{} $\covvaldum$.
To demonstrate this, we state the following simple result showing how spurious \covname{} identification (i.e., \cref{defn:spurious}) relies on the number of queries.
To avoid additional notational burden, we state this result informally, and defer precise notation, more general statements, and proofs to \cref{sec:sample-complexity-proofs}.

\begin{theorem}\label{fact:sufficient-samples-informal}
Suppose that $\dataspace=\Reals$ and there exists $\lipconst>0$ such that all $\model\in\modelspace$ are $\lipconst$-Lipschitz. 
For fixed $\inscale,\outscale>0$ and $j\in[\covdim]$, consider the \downtaskname{} of \cref{defn:spurious} with $\covval=0$.
%
%
%
%
%
For every $n\in\Nats$, there exists a \hyptestname{} $\oraclehyptest\timeind{n}$ that uses only $n$ evaluations of $\model$ yet
\*[
    \specsym(\oraclehyptest\timeind{n}) = 1
    \ \text{ and } \
    \senssym(\oraclehyptest\timeind{n})
    = 1 - \Big(1-\frac{2\outscale}{\lipconst\inscale}\Big)^n.
\]
\end{theorem}

A few remarks on \cref{fact:sufficient-samples-informal}. 
%
First, this result implies that the \username{} can achieve \specname{} \emph{and} \sensname{} arbitrarily close to one by evaluating $\model$ sufficiently many times (i.e., taking $n$ to infinity).
Second, while we did not state this in \cref{sec:impossibility} for simplicity, from the proof of \cref{fact:oracle-hypothesis} it can be observed that \cref{assn:piecewise} can be relaxed to require only Lipschitz piecewise linear functions, and hence this additional structure is not sufficient to circumvent our impossibility result for \completename{} and \additivename{} \fielddefnames{} (i.e., \cref{fact:oracle-hypothesis-examples}). 
Third, we defer the precise definition of the \hyptestname{} to \cref{sec:sample-complexity-proofs}, but here is a simple intuition: uniformly sample $n$ \obsnames{} $\covval\timeind{1:n}$ and reject the \nullname{} if and only if $\model(0,\covval\feat{j}\timeind{t}) > \outscale$ for at least one $t$.
Finally, in \cref{sec:sample-complexity-proofs} we extend this result to the multivariate case and quantify the corresponding dependence on $\covdim$ (unsurprisingly, this dependence is exponential), and we show that this simple \hyptestname{} achieves nearly optimal worst-case dependence on parameters like $\inscale$, $\outscale$, $\lipconst$, $\covdim$, and $n$. 

As a concrete example, consider the spurious \covname{} identification \downtaskname{} for 10 \covnames{} (say, 10 pixels where a watermark appears) on $1$-Lipschitz \modelnames{}.
To identify if the \modelname{} is sensitive (\modelname{} output changes by more than 1\%) to a 5\% change in the \covname{} value, we show it is possible to get perfect \specname{} and over 90\% \sensname{} with roughly 20,000 model evaluations. 
%
%
%
 
While the \downtaskname{} of spurious \covname{} identification can be solved by brute-force evaluations (and a similar argument can be made for other \countmodbehaviourname{} \downtasknames{}), there is much work to be done.
In particular, the simple \hyptestname{} in \cref{fact:sufficient-samples-informal} is designed to always succeed, but may be quite inefficient for more structured \modelnames{}.
Unfortunately, for the existing brute-force \hyptestnames{} that are more efficient, it is unclear for which \downtasknames{} they provably work.
%
%
Our hypothesis testing framework allows us to rigorously evaluate such methods, and our formalization of \downtasknames{} suggests that techniques and methods from optimization theory may be a useful starting point.
 
 

%
%
%
%
%
%
%
%
%
%
%
%
%
%
%


%




%
%
%
%

%
%
%
%


%


%


%
%



%

%



